    \title{Lista de Exercícios: TLS}
        \author{Prof. Gabriel Rodrigues Caldas de Aquino}
        %\date{Compilado em: \\ \today}

\begin{document}

\maketitle

\section{Conceitos: Conexão vs. Sessão}
Um diferencial da arquitetura do TLS é a distinção clara entre uma conexão de transporte e uma sessão de segurança.

\begin{enumerate}
    \item Defina o que é uma \textbf{Conexão} no contexto do TLS.
    \item Defina o que é uma \textbf{Sessão}.
    \item Explique a relação entre Conexões e Sessões. É possível que múltiplas conexões seguras pertençam à mesma sessão?
    \item Qual é o principal benefício de desempenho (performance) ao se utilizar o conceito de Sessão reutilizável? Por que seria custoso negociar novos parâmetros de segurança para cada nova conexão TCP estabelecida?
\end{enumerate}


\section{Arquitetura e Conceitos Fundamentais}
O TLS é projetado para prover segurança entre aplicações, permitindo que se comuniquem de forma segura.

\begin{enumerate}
    \item Diga qual a localização do TLS na pilha de protocolos TCP/IP. Justifique.
    \item Do ponto de vista do desenvolvedor da aplicação, como o TLS é percebido em comparação ao TCP padrão?
    \item O TLS é composto por um conjunto de protocolos. Descreva o \textbf{Protocolo de Registro (Record Protocol)} e a sua função.
    \item Quais são as duas camadas do TLS?
    \item Cite outros três protocolos de gerenciamento que operam "em cima" do \textbf{Protocolo de Registro} e diga qual a função de cada um deles.
\end{enumerate}

\section{O Protocolo de Handshake}
O Handshake é a parte mais complexa do TLS, responsável pela negociação de parâmetros e autenticação antes da troca de dados.

\begin{enumerate}
    \item Descreva as Quatro Fases do Handshake.
    \item O Handshake ocorre em diferentes fases. Descreva o que é negociado na \textbf{Fase 1} através das mensagens \texttt{client\_hello} e \texttt{server\_hello}.
    \item Na Fase 1, tanto o cliente quanto o servidor geram um valor \textbf{Random}. Qual é a importância desse valor para a segurança da conexão ?
    \item Explique a função do campo \textbf{CipherSuite}.
    \item Dado que um cliente se conecta ao servidor, diga qual o motivo de ocorrer a autenticação do servidor durante o handshake.
\end{enumerate}


\section{Outros protocolos do TLS}
O TLS também possui protocolos para manutenção e sinalização.

\begin{enumerate}
    \item Qual a finalidade do \textbf{Protocolo de Alerta (Alert Protocol)}? Qual é a estrutura da mensagem?
    \item O que acontece com a conexão TLS quando \textbf{Protocolo de Alerta (Alert Protocol)} apresenta uma mensagem de alerta de nível \texttt{fatal}? Dê um exemplo de erro que geraria tal alerta.
    \item Diga qual a finalidade do Protocolo \textbf{Change Cipher Spec}. O que essa mensagem causa?
\end{enumerate}



\section{O Heartbeat}
O \textbf{Heartbeat} (RFC 6250) possui dois propósitos principais.
\begin{enumerate}
    \item Explique o propósito de "Keep-Alive".
    \item Explique o propósito de "Firewall Traversal" (Evitar Ociosidade).
\end{enumerate}

O \textbf{Heartbleed (2014)}  foi causado por um erro de programação na biblioteca OpenSSL ao implementar a extensão Heartbeat.

\begin{enumerate}
    \item Descreva o comportamento esperado do protocolo Heartbeat (o que o servidor deveria fazer ao receber um request).
    \item Descreva a falha na vulnerabilidade \textbf{Heartbleed (2014)}
    \item Explique a falha lógica: O que o servidor OpenSSL deixava de verificar ao processar a mensagem \texttt{heartbeat\_request}?
    \item Como essa falha permitia que um atacante podesse ler dados sensíveis (como chaves privadas e senhas) da memória do servidor?
\end{enumerate}
\end{document}